\chapter{Ход работы}
\label{ch:practice}

В этой главе по порядку описаны все четыре урока методички:
консольные задачи на базовые конструкции JavaScript (уроки~1 и~2),
реализация сайта-каталога (урок~3) и публикация сайта на собственном
сервере с обратным прокси (урок~4). Для каждого урока приведён код
решения, ожидаемый вывод и скриншоты результата.

\section{Урок~1. Базовые понятия}

Цель урока --- научиться объявлять переменные и выводить значения в
консоль. Согласно методичке, моделируется бонусная программа
интернет-магазина.

\subsection{Задача 1: имя и баланс пользователя}

Объявляются две переменные --- \texttt{username} и
\texttt{bonusBalance}, --- значения которых выводятся в консоль двумя
строками.

\begin{lstlisting}[caption={\texttt{les-1/task1.js} --- имя и баланс},captionpos=b]
const username = 'my name';
const bonusBalance = 1000;

console.log('Пользователь ' + username);
console.log('Баланс ' + bonusBalance);
\end{lstlisting}

Скрипт запускается командой \texttt{node task1.js} из каталога
\texttt{les-1/}. Вывод:

\begin{lstlisting}[caption={Вывод \texttt{task1.js}},captionpos=b]
Пользователь my name
Баланс 1000
\end{lstlisting}

% — Скриншот: вывод les-1/task1 —
\begin{figure}[h]
  \centering
  \includegraphics[width=0.85\textwidth]{img/01_les1_task1}
  \caption{Вывод задачи 1 урока~1 в консоли Node.js}
  \label{fig:les1_task1}
\end{figure}

\subsection{Задача 2: баланс через 7 дней}

За каждую покупку начисляется \(50\) бонусных баллов, ежедневно
сгорает \(3\) балла. Покупки совершаются раз в два дня, первая ---
во вторник. За 7~дней приходится 4~покупки (вт, чт, сб, пн), баланс
изменяется на \(4 \cdot 50 - 7 \cdot 3 = 200 - 21 = 179\) баллов.

\begin{lstlisting}[caption={\texttt{les-1/task2.js} --- баланс через 7 дней},captionpos=b]
let bonusBalance = 1000;
const purchaseBonus = 50;
const dailyBurn = 3;
const days = 7;

// Первая покупка во вторник, дальше раз в два дня: вт, чт, сб, пн.
const purchases = 4;

bonusBalance = bonusBalance + purchases * purchaseBonus - days * dailyBurn;

console.log('Баланс через ' + days + ' дней: ' + bonusBalance);
\end{lstlisting}

Вывод:

\begin{lstlisting}
Баланс через 7 дней: 1179
\end{lstlisting}

% — Скриншот: вывод les-1/task2 —
\begin{figure}[h]
  \centering
  \includegraphics[width=0.85\textwidth]{img/02_les1_task2}
  \caption{Вывод задачи 2 урока~1 в консоли Node.js}
  \label{fig:les1_task2}
\end{figure}

\section{Урок~2. Основы JavaScript}

Цель урока --- работа с массивами, циклами и строковыми методами на
примере мини-мессенджера.

\subsection{Задача 1: вывод переписки в чате}

Инициализируется массив из четырёх сообщений; чётные индексы
(\(0,~2,~\dots\)) принадлежат Другу, нечётные --- Вам. Используется
цикл \texttt{for} с условным оператором, выбирающим имя отправителя.

\begin{lstlisting}[caption={\texttt{les-2/task1.js} --- вывод переписки},captionpos=b]
const messages = [
    'Пойдем гулять в парк?',
    'Кажется, дождь собирается. Лучше пойдем в кино!',
    'Давай, сегодня как раз вышел новый фильм.',
    'Встречаемся через час у кинотеатра.',
];

// Первое сообщение — от Друга, дальше чередование.
for (let i = 0; i < messages.length; i++) {
    const author = i % 2 === 0 ? 'Друг' : 'Вы';
    console.log(author + ': ' + messages[i]);
}
\end{lstlisting}

Вывод:

\begin{lstlisting}
Друг: Пойдем гулять в парк?
Вы: Кажется, дождь собирается. Лучше пойдем в кино!
Друг: Давай, сегодня как раз вышел новый фильм.
Вы: Встречаемся через час у кинотеатра.
\end{lstlisting}

% — Скриншот: вывод les-2/task1 —
\begin{figure}[h]
  \centering
  \includegraphics[width=0.85\textwidth]{img/03_les2_task1}
  \caption{Вывод переписки в консоли Node.js}
  \label{fig:les2_task1}
\end{figure}

\subsection{Задача 2: поиск по тексту сообщений}

Реализован простейший поиск: тот же массив сообщений фильтруется
методом \texttt{Array.prototype.filter} с проверкой
\texttt{String.prototype.includes(searchText)}. Найденные сообщения
выводятся списком.

\begin{lstlisting}[caption={\texttt{les-2/task2.js} --- поиск через \texttt{includes}},captionpos=b]
const messages = [
    'Пойдем гулять в парк?',
    'Кажется, дождь собирается. Лучше пойдем в кино!',
    'Давай, сегодня как раз вышел новый фильм.',
    'Встречаемся через час у кинотеатра.',
];

const searchText = 'кино';

console.log('Поиск по слову: "' + searchText + '"');
const found = messages.filter(msg => msg.includes(searchText));

if (found.length === 0) {
    console.log('Совпадений не найдено.');
} else {
    found.forEach(msg => console.log('— ' + msg));
}
\end{lstlisting}

Вывод:

\begin{lstlisting}
Поиск по слову: "кино"
— Кажется, дождь собирается. Лучше пойдем в кино!
— Встречаемся через час у кинотеатра.
\end{lstlisting}

% — Скриншот: вывод les-2/task2 —
\begin{figure}[h]
  \centering
  \includegraphics[width=0.85\textwidth]{img/04_les2_task2}
  \caption{Вывод поиска по сообщениям в консоли Node.js}
  \label{fig:les2_task2}
\end{figure}

\section{Урок~3. Сайт <<Персонажи Marvel>>}

Главный артефакт лабораторной работы --- одностраничный сайт-каталог
девятнадцати героев Marvel, реализованный в файлах
\texttt{les-3/index.html}, \texttt{les-3/index.js},
\texttt{les-3/start.js}.

\subsection{Каркас страницы}

\texttt{index.html} --- статическая разметка на компонентах Bootstrap.
Сразу после загрузки страница содержит тёмный навбар, заголовок
<<Персонажи Marvel>> и спиннер-плейсхолдер; ниже --- два пустых
контейнера для будущих карточек и модальных окон.

\begin{lstlisting}[caption={Ключевая часть \texttt{les-3/index.html}},captionpos=b]
<div class="container">
    <h1 class="pt-3 pb-3">Персонажи Marvel</h1>

    <div class="row" id="character-card-box">
        <div class="d-flex justify-content-center">
            <div class="spinner-border text-danger" role="status">
                <span class="visually-hidden">Loading...</span>
            </div>
        </div>
    </div>

    <div id="character-modal-box"></div>
</div>

<script src=".../bootstrap.bundle.min.js"></script>
<script src="index.js"></script>
<script src="start.js"></script>
<script>
    document.addEventListener("DOMContentLoaded", function() { start(); });
</script>
\end{lstlisting}

\subsection{Получение данных с API}

Функция \texttt{fetchCharacters} выполняет GET-запрос к учебному
API и возвращает распарсенный JSON-массив. При ошибке HTTP
выбрасывается исключение, которое перехватывается в \texttt{start()}.

\begin{lstlisting}[caption={\texttt{les-3/index.js} --- запрос к API},captionpos=b]
const API_URL = 'https://jsfree-les-3-api.onrender.com/characters';

async function fetchCharacters() {
    const response = await fetch(API_URL);
    if (!response.ok) {
        throw new Error(`HTTP ${response.status}`);
    }
    return await response.json();
}
\end{lstlisting}

\subsection{Карточка персонажа}

\texttt{getCharacterCard} принимает один объект-персонаж и возвращает
HTML-строку с Bootstrap-карточкой: обложка (\texttt{thumbnail}),
имя по-русски, имя на оригинале, краткое описание (обрезается до
\(160\) символов) и кнопка <<Подробнее>>, которая открывает модальное
окно через атрибуты \texttt{data-bs-toggle} и \texttt{data-bs-target}.

\begin{lstlisting}[caption={\texttt{getCharacterCard} --- HTML-карточка героя},captionpos=b]
function getCharacterCard(character) {
    const modalId = `character-modal-${character.id}`;
    const thumb = escapeHtml(toHttps(character.thumbnail));
    const name = escapeHtml(character.name);
    const nameor = escapeHtml(character.nameor);
    const short = escapeHtml(truncate(character.description, 160));

    return `
<div class="col-12 col-md-6 col-lg-4 mb-4">
    <div class="card h-100 shadow-sm">
        <img src="${thumb}" class="card-img-top" alt="${name}" loading="lazy"
             style="height: 320px; object-fit: cover; object-position: top;">
        <div class="card-body d-flex flex-column">
            <h5 class="card-title mb-1">${name}</h5>
            <p class="text-muted small mb-2">${nameor}</p>
            <p class="card-text flex-grow-1">${short}</p>
            <button type="button" class="btn btn-danger mt-2"
                    data-bs-toggle="modal" data-bs-target="#${modalId}">
                Подробнее
            </button>
        </div>
    </div>
</div>`;
}

function getCharacterCards(characters) {
    return characters.map(getCharacterCard).join('');
}
\end{lstlisting}

Все значения, попадающие в HTML, пропускаются через
\texttt{escapeHtml} --- защита от XSS на случай, если в данных
встретятся служебные символы (\texttt{<}, \texttt{>}, \texttt{\&},
кавычки). URL обложки переписывается с \texttt{http} на
\texttt{https} (функция \texttt{toHttps}), иначе браузер заблокирует
изображение из-за \textit{mixed content} на HTTPS-странице.

\subsection{Модальное окно персонажа}

\texttt{getCharacterModal} собирает Bootstrap-модал с подробной
биографией героя, увеличенной обложкой и списком до \(10\) комиксов с
его участием. Идентификатор модала формируется из \texttt{id} героя ---
именно по нему карточка ищет нужный модал через
\texttt{data-bs-target}.

\begin{lstlisting}[caption={\texttt{getCharacterModal} --- модальное окно},captionpos=b]
function getCharacterModal(character) {
    const modalId = `character-modal-${character.id}`;
    const comics = Array.isArray(character.comics) ? character.comics.slice(0, 10) : [];
    const comicsList = comics.length
        ? `<ul>${comics.map(c => `<li>${escapeHtml(c.name)}</li>`).join('')}</ul>`
        : '<p class="text-muted">Список комиксов недоступен.</p>';

    return `
<div class="modal fade" id="${modalId}" tabindex="-1">
    <div class="modal-dialog modal-lg modal-dialog-scrollable modal-dialog-centered">
        <div class="modal-content">
            <div class="modal-header"> ... </div>
            <div class="modal-body">
                <div class="row">
                    <div class="col-md-4"> <img src="..." class="img-fluid rounded"> </div>
                    <div class="col-md-8">
                        <p>${escapeHtml(character.description)}</p>
                        <h6>Появления в комиксах:</h6>
                        ${comicsList}
                    </div>
                </div>
            </div>
            <div class="modal-footer"> ... </div>
        </div>
    </div>
</div>`;
}

function getCharacterModals(characters) {
    return characters.map(getCharacterModal).join('');
}
\end{lstlisting}

(В листинге опущены повторяющиеся блоки \texttt{modal-header} и
\texttt{modal-footer} --- полный код находится в репозитории.)

\subsection{Оркестрация: \texttt{start.js}}

Функция \texttt{start} вызывается по событию
\texttt{DOMContentLoaded}: получает данные с API, заполняет оба
контейнера готовой разметкой; при ошибке заменяет спиннер сообщением
для пользователя.

\begin{lstlisting}[caption={\texttt{les-3/start.js} --- запуск приложения},captionpos=b]
async function start() {
    const cardBox = document.getElementById('character-card-box');
    const modalBox = document.getElementById('character-modal-box');

    try {
        const characters = await fetchCharacters();
        cardBox.innerHTML = getCharacterCards(characters);
        modalBox.innerHTML = getCharacterModals(characters);
    } catch (err) {
        cardBox.innerHTML = `
<div class="col-12">
    <div class="alert alert-danger">
        Не удалось загрузить персонажей: ${err.message}.
    </div>
</div>`;
    }
}
\end{lstlisting}

\subsection{Локальная проверка}

Открыть \texttt{index.html} напрямую через \texttt{file://} нельзя:
браузер блокирует \texttt{fetch} к чужому домену. Поэтому для проверки
поднимается локальный HTTP-сервер:

\begin{lstlisting}
cd les-3 && python3 -m http.server 8000
\end{lstlisting}

Страница открывается по адресу \texttt{http://localhost:8000}. Через
несколько секунд (или до \(30\) секунд при холодном старте API)
спиннер сменяется сеткой карточек. Клик по кнопке <<Подробнее>>
открывает модальное окно с биографией и списком комиксов.

% — Скриншот: главная страница локально —
\begin{figure}[h]
  \centering
  \includegraphics[width=0.85\textwidth]{img/05_marvel_local_grid}
  \caption{Сайт <<Персонажи Marvel>> локально: сетка карточек}
  \label{fig:marvel_local}
\end{figure}

% — Скриншот: открытая модалка —
\begin{figure}[h]
  \centering
  \includegraphics[width=0.85\textwidth]{img/06_marvel_modal}
  \caption{Открытое модальное окно с биографией и комиксами}
  \label{fig:marvel_modal}
\end{figure}

% — Скриншот: мобильная версия —
\begin{figure}[h]
  \centering
  \includegraphics[width=0.45\textwidth]{img/07_marvel_mobile}
  \caption{Адаптивная вёрстка на мобильной ширине (\(360\,\)px)}
  \label{fig:marvel_mobile}
\end{figure}

\section{Урок~4. Публикация сайта}

Согласно методичке, готовый сайт должен быть выложен на публичный
хостинг, а ссылка указана в отчёте. Вместо коммерческого Timeweb
сайт развёрнут на собственном домашнем сервере, на котором уже
работает обратный прокси \textbf{Caddy} с автоматическими
TLS-сертификатами Let's~Encrypt (см.~\hyperref[ch:project]{главу~2}).

\subsection{Автоматизация деплоя}

В корне \texttt{lab13/} лежит скрипт \texttt{deploy.sh} с двумя
ключами --- \texttt{--install} и \texttt{--uninstall}. Оба режима
идемпотентные, перезапуск ничего не ломает.

\begin{lstlisting}[caption={Запуск деплоя и сноса сайта},captionpos=b]
./deploy.sh --install     # развернуть сайт
./deploy.sh --uninstall   # снести сайт после показа
\end{lstlisting}

Шаги режима \texttt{--install}:

\begin{enumerate}
  \item Запрашивает у пользователя домен (по умолчанию ---
        \texttt{marvel.server34.netcraze.club}) и путь к исходникам
        сайта (по умолчанию --- соседний каталог
        \texttt{les-3/}).
  \item Проверяет, что Caddy уже установлен
        (\texttt{50-install-caddy-proxy.sh} был запущен ранее);
        иначе --- ошибка и выход.
  \item Делает резервные копии \texttt{docker-compose.yml} и
        \texttt{Caddyfile} с меткой времени.
  \item Через \texttt{rsync -a --delete} синхронизирует исходники
        \texttt{lab13/les-3/} в \texttt{/opt/stack/marvel/site/}.
  \item Если в \texttt{docker-compose.yml} ещё нет сервиса
        \texttt{marvel} --- добавляет: образ \texttt{nginx:alpine},
        монтирует директорию сайта в \texttt{/usr/share/nginx/html}
        в режиме только для чтения.
  \item Если в \texttt{Caddyfile} ещё нет блока для нашего поддомена
        --- дописывает в конец:
        \begin{lstlisting}
marvel.server34.netcraze.club {
    encode gzip
    reverse_proxy marvel:80
}
        \end{lstlisting}
  \item Валидирует итоговый \texttt{docker-compose.yml}
        (\texttt{docker compose config}), запускает контейнер
        \texttt{marvel} и перезагружает Caddy
        (\texttt{caddy reload}).
\end{enumerate}

Режим \texttt{--uninstall} выполняет обратные действия: останавливает
и удаляет контейнер \texttt{marvel}, awk-фильтром вырезает сервис из
\texttt{docker-compose.yml} и блок поддомена из \texttt{Caddyfile},
удаляет каталог \texttt{/opt/stack/marvel/} и перезагружает Caddy.
Бэкапы конфигов с меткой времени остаются рядом --- на случай отката.

Первый запрос к новому поддомену запускает выпуск TLS-сертификата
Let's~Encrypt; через несколько секунд страница открывается по HTTPS.
Готовый адрес --- \url{https://marvel.server34.netcraze.club}.

% — Скриншот: живой сайт по HTTPS —
\begin{figure}[h]
  \centering
  \includegraphics[width=0.85\textwidth]{img/08_marvel_live_https}
  \caption{Развёрнутый сайт по адресу
           \texttt{https://marvel.server34.netcraze.club}
           (Let's~Encrypt-сертификат активен)}
  \label{fig:marvel_live}
\end{figure}

% — Скриншот: вывод docker compose ps на сервере —
\begin{figure}[h]
  \centering
  \includegraphics[width=0.85\textwidth]{img/09_docker_compose_ps}
  \caption{Сервис \texttt{marvel} в общем \texttt{docker-compose.yml}
           домашнего сервера}
  \label{fig:compose_ps}
\end{figure}

\endinput
